\section{Results and Discussion}
\label{sec:results_discussion}

\subsection{Model Performance Results}
The predictive performance of the Logistic Regression and Random Forest models was evaluated using several metrics on a 20\% hold-out test set. The optimized Random Forest model demonstrated superior performance across all primary metrics. Table~\ref{tab:model_comparison} summarizes these results.

\begin{table}[h]
\centering
\caption{Comparison of Model Performance Metrics}
\label{tab:model_comparison}
\begin{tabular}{lcc}
\hline
\textbf{Metric} & \textbf{Logistic Regression} & \textbf{Random Forest} \\ \hline
Accuracy         & $\approx$ 0.932             & 0.934                  \\
Precision        & $\approx$ 0.958             & 0.963                  \\
Recall           & $\approx$ 0.938             & 0.942                  \\
F1-score         & $\approx$ 0.941             & 0.952                  \\
ROC-AUC          & $\approx$ 0.985             & 0.988                  \\ \hline
\end{tabular}
\end{table}

\subsubsection{Statistical Significance}
A paired t-test was performed comparing the F1-scores of both models over 10-fold cross-validation. The test yielded a p-value of 0.0134 ($p < 0.05$), allowing us to reject the null hypothesis. This confirms that the Random Forest model's improvement over Logistic Regression is statistically significant ($t = -3.066$).

\subsection{Feature Importance Analysis}
The Random Forest model's feature importance analysis identified the most influential behavioral predictors of smartphone addiction. The results highlight a heavy reliance on usage duration and psychological state.

\begin{itemize}
    \item \textbf{Daily Screen Time Hours (53.0\%):} The most significant predictor, indicating a direct relationship between total device usage and addiction risk.
    \item \textbf{Social Media Hours (24.4\%):} The second most critical factor, suggesting that the type of content (specifically social networking) is a major driver.
    \item \textbf{Stress Level (8.5\%):} Indicates a correlation between higher reported stress and addictive usage patterns.
    \item \textbf{Secondary Factors:} Sleep hours, work/study hours, and notification frequency each contributed approximately 2\% to the model's predictive power.
\end{itemize}

\subsection{Discussion}
The findings suggest that smartphone addiction is primarily driven by the quantity of time spent on the device, particularly on social media platforms. The high ROC-AUC (0.988) indicates that the Random Forest model is exceptionally capable of distinguishing between ``Addicted'' and ``Not Addicted'' individuals.

The statistically significant superiority of the Random Forest model suggests that non-linear interactions between variables—such as the interplay between stress levels and screen time—play a crucial role in predicting addiction. While linear models like Logistic Regression perform well, they fail to capture the complexity of behavioral patterns as effectively as ensemble-based methods.

Furthermore, the relatively low importance of the engineered \textit{screen\_time\_to\_social\_network\_ratio} (0.3\%) compared to the raw \textit{social\_media\_hours} suggests that the absolute volume of social media usage is a more direct indicator of addiction than its proportion relative to other activities.
